\appendix
\section{Prompts and tool definitions}
\label{app:prompts}

\paragraph{System prompt in the \texttt{self} and placebo conditions.} This text follows Qwen's own tool block, which lists the \texttt{set\_mode} JSON schema:
\begin{quote}\small\ttfamily
You can change how your own neurons are set before you answer, by calling set\_mode. Each mode is a small edit to your MLP layers. Your modes:\\
- lazy: Short, direct answers with no padding. Good for quick facts, definitions, code snippets, one-line commands and simple arithmetic. Avoid it for multi-step planning or scheduling puzzles, where it tends to take greedy shortcuts.\\
- normal: no edit.\\
Before answering, decide whether a mode clearly fits the request. If one does, call set\_mode once, then answer after the tool result. If none fits, just answer normally without calling any tool.
\end{quote}
With a controller connected, one more line is added: ``If the user wants a style none of your modes gives, you may call request\_mode to ask Claude to build one, then switch it on with set\_mode.''

\paragraph{Tool result after \texttt{set\_mode}.} ``Mode `lazy' is on for your next answer. Now answer the user.'' The placebo condition returns the same text but arms no edit.

\paragraph{The \texttt{lazy} description and the planning warning.} The warning against planning and scheduling comes from Lazy Genius test~2, a different question set from the test-1 set used in E4.

\paragraph{Router brief (\texttt{routed}).} A Claude subagent was told the goal (``shorter answers without hurting correctness''), given the mode description above, and allowed to read only a file of the 96 question texts, with no rubrics, answers or grades. It labeled 65 questions lazy and 31 normal.

\paragraph{Grader brief.} Each grader read only Lazy Genius's \texttt{GRADING\_SAMPLED.md} and one blind chunk. It could run tiny isolated code fixtures, and it returned \texttt{core\_correct}, \texttt{correct}, \texttt{completion\_failure} and a reason for each item. The Lazy Genius tests ran graders at a ``high'' effort setting. Here they ran at the default setting, which is a deviation. The calibration items (\S\ref{sec:e4}) measure its effect.

\section{Deviations from the pre-registration}
\begin{itemize}
\item \textbf{Placebo condition (E4).} It was added after the E4 results and is exploratory.
\item \textbf{Grader effort setting.} The default was used instead of ``high''; see above.
\item \textbf{Decode benchmark (E2).} It used greedy decoding, so each edit produced different text. Speeds are per token.
\end{itemize}

\section{Reproducing}
The code, presets, raw generations, blind keys, judgments and analysis outputs are in the repository. All paper numbers are generated by \texttt{paper/make\_numbers.py} from \texttt{results/*.json}.
